% 第一部分：算法性质与当前实现的边界分别说明。
\begingroup
\setlength{\parskip}{2bp}
\setlength{\textfloatsep}{8bp plus 2bp minus 2bp}
\setlength{\intextsep}{8bp plus 2bp minus 2bp}
\setcounter{figure}{0}
\renewcommand{\thefigure}{1-\arabic{figure}}
\renewcommand{\theHfigure}{section1.\arabic{figure}}
\renewcommand{\figurename}{图}
\captionsetup{font=small,labelsep=quad,skip=5bp}
\newcommand{\OverviewTitle}[2]{%
  \par\Needspace{4\baselineskip}%
  #1{#2}{\bfseries #2}\par\vspace{2bp}}
\newcommand{\OverviewFigure}[4][0.90]{%
  \begin{figure}[htbp]
    \centering
    \includegraphics[width=#1\linewidth]{section1Figures/#2.pdf}
    \caption{#3}\label{#4}
  \end{figure}}

\TemplateSection{1.\quad 启发式算法综述}
{\bfseries\fontsize{12bp}{15.6bp}\selectfont 1.\quad 启发式算法综述}\par\vspace{6bp}

N-Puzzle 的求解可以看作从初始棋盘到目标棋盘的路径搜索。随着棋盘规模增大，
直接枚举状态的时间和存储代价都会迅速增长。启发式搜索利用问题自身的结构，
估计当前状态与目标之间尚有多远，使搜索优先考察更有希望形成低代价解的路径。
本部分先说明 A* 的评价规则及其最优性条件，再分析重复状态与启发函数的一致性，
最后介绍用迭代加深控制存储开销的 IDA*，为后续实现与实验比较提供依据。

\OverviewTitle{\TemplateSubsection}{1.1 算法描述}

设状态图为 $G=(V,E)$，初态为 $S$，目标状态为 $G_{\mathrm{goal}}$，
边代价 $c(s,s')$ 表示从状态 $s$ 转移到 $s'$ 的实际消耗。
本实验的一次合法移动是空位与相邻将牌交换，代价为 1，因此解的总代价就是移动步数。
为避免把问题状态与搜索过程混淆，本文用 $s$ 表示棋盘状态，用 $n$ 表示搜索结点。
结点除保存状态外，还保存父结点、到达动作及路径代价；同一棋盘经两条不同路径到达，
可以产生两个不同的搜索结点。

A* 使用下面的评价函数决定扩展顺序：
\begin{equation}
  f(n)=g(n)+h(n).
  \label{eq:overview-f}
\end{equation}
其中，$g(n)$ 是沿当前父链从初态到结点 $n$ 已经付出的代价，能够由路径直接计算；
$h(n)$ 是从该结点状态到目标的剩余代价估计；$f(n)$ 则估计经过当前路径完成求解的总代价。
真实最小剩余代价记为 $h^*(s)$。由于直接计算 $h^*$ 往往等价于求解原问题，
实际算法采用更容易计算的下界作为启发信息。

这一规则同时考虑已付代价与剩余距离。令 $h=0$ 时，A* 退化为按累计代价展开的一致代价搜索；
只按 $h$ 排序则是贪心最佳优先搜索，不能直接沿用 A* 的最优性结论。
对 N-Puzzle 而言，不在位将牌数、曼哈顿距离和满足成本分摊条件的模式数据库都可提供启发估计，
但其计算成本与估计强度不同。更精确的下界可能减少搜索结点，也可能增加单次估值或预处理的开销，
因此算法效率仍需结合实际测量判断。

\OverviewTitle{\TemplateSubsection}{1.2 算法的可采纳性}

启发函数的可采纳性指它不高估剩余最优代价，即对能够到达目标的状态满足
\begin{equation}
  0\leq h(s)\leq h^*(s),\qquad h(G_{\mathrm{goal}})=0.
  \label{eq:overview-admissible}
\end{equation}
这是一项关于估值的条件。搜索算法的可采纳性则指：在给定搜索模型和终止条件下，
若存在解，算法能够返回代价最小的解。两者并不等同，后者还取决于搜索顺序、重复状态处理
以及是否完整保留可能通向最优解的路径。下文均以有限状态图、正边代价及可达目标为讨论前提，
本实验的单位代价棋盘满足这些条件。

\OverviewTitle{\TemplateSubsubsection}{1.2.1 A*算法}

A* 维护 Open 表和 Closed 表。Open 保存已经生成、尚待处理的候选结点，
Closed 保存已经完成扩展的状态。每一轮从 Open 取出 $f$ 最小的结点，
若它表示目标状态，就沿父链恢复解路径；否则生成后继并维护候选集合。
图~\ref{fig:overview-astar} 给出了本项目使用的不重开 Closed 的处理流程。

具体步骤为：（1）清空两张表，令根结点 $g(S)=0$、$f(S)=h(S)$，将其加入 Open；
（2）若 Open 为空，搜索结束并报告未找到解，否则取出最小 $f$ 的结点；
（3）先检查该结点是否为目标，若是则返回路径；
（4）将非目标结点的状态加入 Closed，生成合法后继，并令 $g'=g+c$；
（5）对尚未进入 Closed 的后继，若其状态不在 Open 则插入，若已经存在则保留 $g$ 更小的结点，
随后回到步骤（2）。同一状态的 $h$ 相同，所以比较 $g$ 与比较 $f$ 等价。

\OverviewFigure[0.82]{fig-1-1-astar-flow}
  {A* 图搜索流程与重复状态处理}
  {fig:overview-astar}

目标判定必须结合出队顺序理解。生成目标只说明发现了一条可行路径，
Open 中仍可能存在通向更短解的候选；只有目标作为最小 $f$ 结点被取出时，
才能在相应理论条件下确认其最优性。本项目的教师类 \texttt{BestFirstSearcher}
正是在取出结点之后判定目标，并使用学生实现的 \texttt{PqFrontier} 管理 Open。
该 Open 由优先队列与状态索引表共同维护，Closed 则使用状态集合，
因此“选出最小估值”和“判断是否为同一棋盘”是两个不同的操作。
\FloatBarrier

\OverviewTitle{\TemplateSubsubsection}{1.2.2 A*算法可采纳性证明}

A* 的经典论证利用最优路径上的边界结点\cite{hart1968}。
但需要先说明边界结点为什么不会因重复状态处理而丢失：
树搜索保留不同路径；允许重开的图搜索在发现更小 $g$ 时重新处理该状态；
本项目的图搜索则直接跳过 Closed 中的状态。因此，对当前实现，
除可采纳性外还要求启发函数满足一致性：
\begin{equation}
  h(s)\leq c(s,s')+h(s')\qquad \text{对每条合法边 }(s,s').
  \label{eq:overview-consistency}
\end{equation}

\textbf{首先证明：在一致性条件下，状态第一次出队时的 $g$ 已经最优。}
记从初态到状态 $x$ 的最短代价为 $g^*(x)$。反设 $x$ 是第一个以非最优代价出队的状态，
即 $g(x)>g^*(x)$。沿一条到 $x$ 的最短路径，取当时第一个尚未进入 Closed 的状态 $p$。
它若为初态，最优前缀由根结点表示；否则其前驱已以最优代价扩展，
因此 $p$ 的最优前缀已经生成，且 Open 的替换规则会保留这条最优前缀。
将式~\eqref{eq:overview-consistency} 沿 $p$ 到 $x$ 的子路径累加，可得
\begin{equation}
  f(p)=g^*(p)+h(p)
  \leq g^*(x)+h(x)
  <g(x)+h(x)=f(x).
  \label{eq:overview-pop-optimal}
\end{equation}
这与 $x$ 是 Open 中最小 $f$ 的出队结点矛盾。因此结论成立，Closed 中的状态无需重开。

\textbf{再证明：被取出的第一个目标结点具有最小解代价。}
设最优解代价为 $C^*$。在最优目标尚未被取出时，最优路径上仍有边界结点 $n$ 位于 Open，
且其最优前缀没有丢失。根据可采纳性，
\begin{equation}
  f(n)=g^*(n)+h(n)\leq g^*(n)+h^*(n)=C^*.
  \label{eq:overview-frontier}
\end{equation}
若算法先取出一个代价为 $C>C^*$ 的目标结点 $G'$，由于目标的启发值为 0，
必有 $f(G')=C>C^*\geq f(n)$，与最小 $f$ 出队规则矛盾。
于是第一个出队目标的代价只能是 $C^*$。

\OverviewFigure{fig-1-2-optimality-frontier}
  {最优路径的边界结点与候选目标的估值关系}
  {fig:overview-optimality}

图~\ref{fig:overview-optimality} 用一组示意值展示了关键不等式：
最优路径总代价为 6，边界结点已付代价为 2，启发估计为 3，因此 $f(n)=5$；
另一条路径虽已生成代价为 8 的目标结点，也不能越过该边界结点先出队。
图中的 $G^*$ 与 $G'$ 是通过不同路径到达\textbf{同一目标状态}的两个搜索结点，
并非改变了棋盘的目标。该图用于解释证明，不代表某次实验的结点记录。

若启发函数仅可采纳而不一致，上述“第一次出队即有最优 $g$”的前提可能不成立，
此时需要允许更短路径触发重开等相应机制。不能只检查 $h\leq h^*$，
就断言当前任何不重开的图搜索都能保证最优；同样，改变为带权估值后也需要重新分析。
\FloatBarrier

\OverviewTitle{\TemplateSubsection}{1.3 重复扩展结点的问题}

状态空间通常是一张图，而不是天然没有重复的树。在 N-Puzzle 中，
空位向左移动后再向右移动会回到原棋盘，多条更长的动作序列也可能到达同一排列。
若将每条路径都独立展开，算法会反复求解相同的后续问题。
但简单地“见过就删除”也可能丢弃更短路径，所以去重必须同时考虑状态身份和当前路径代价。

\OverviewTitle{\TemplateSubsubsection}{1.3.1 重复扩展结点问题}

应当区分重复生成与重复扩展。重复生成是不同路径产生了相同状态，
即使使用 Closed 也不一定能完全避免；重复扩展则意味着该状态的后继被再次枚举和估值。
图~\ref{fig:overview-duplicate} 是一个一般的抽象状态图，用来说明两条路径对同一状态的竞争，
不将其当作实际的 N-Puzzle 棋盘转换实例。

\OverviewFigure{fig-1-3-duplicate-state}
  {同一状态经不同路径到达时的代价与父链更新}
  {fig:overview-duplicate}

假设旧路径 $S\to B\to C\to x$ 先被发现，其 $g(x)=3$；
随后发现新路径 $S\to A\to x$，其 $g(x)=2$。如果状态 $x$ 尚在 Open，
就应保留新结点，并使它的父链指向 $A$，否则恢复出的路径仍可能沿用旧路线。
如果 $x$ 已在 Closed，则必须依据前述正确性条件决定是否重开；
在一致启发函数下，这种“已经以较大 $g$ 关闭、后来又出现较小 $g$”的情况不会发生。

本项目用 \texttt{PriorityQueue<Node>} 按评价值排序，用
\texttt{HashMap<State, Node>} 定位同一状态在 Open 中的有效结点，
Closed 则由 \texttt{HashSet<State>} 保存。普通入队和取出最小值的堆操作为
$O(\log N)$，状态索引的查询平均为 $O(1)$；但替换旧结点时调用的
\texttt{PriorityQueue.remove(oldNode)} 需要线性查找，复杂度为 $O(N)$。
因此，加入哈希索引解决了状态定位问题，却没有使所有替换操作都变成对数复杂度。
实际收益取决于候选规模及较优重复路径出现的频率。

哈希值还不能代替状态相等。不同棋盘可能具有相同哈希，
查重时必须由完整状态比较确认是否相同；作为键的棋盘内容也应保持不变，
以免插入后无法正确查回。对 IDA*，当前路径集合只剪去形成环的后继，
不会把以前阈值下访问过的所有状态永久关闭，原因将在 1.4 中说明。
\FloatBarrier

\OverviewTitle{\TemplateSubsubsection}{1.3.2 启发函数的单调性}

式~\eqref{eq:overview-consistency} 也称单调性条件，其含义是：
一步转移带来的启发估值下降，不能超过这一步的实际代价。对父子结点 $n,n'$，
利用 $g(n')=g(n)+c(n,n')$，可得到
\begin{equation}
  f(n')=g(n)+c(n,n')+h(n')
  \geq g(n)+h(n)=f(n).
  \label{eq:overview-monotone}
\end{equation}
因此一致启发函数使同一路径上的 $f$ 不下降。
这与可采纳性对“估计值和真实剩余代价”的比较不同，是对相邻状态估值关系的约束。

\OverviewFigure[0.80]{fig-1-4-consistency}
  {一致性与可采纳性的区别：同一两步路径上的两种估值}
  {fig:overview-consistency}

图~\ref{fig:overview-consistency} 中，两条示意路径的实际剩余代价依次都是 2、1、0。
上图取 $h=(2,1,0)$，相邻状态均满足一致性，故 $f=(2,2,2)$。
下图取 $h=(2,0,0)$，每个估计仍不超过真实剩余代价，所以可采纳；
但首条边出现 $2>1+0$，使 $f$ 从 2 降到 1，因此不一致。
这说明可采纳性并不能反推出单调性。

反过来，若目标启发值为 0，一致性可以推出可采纳性。
沿任意到目标的路径 $s=s_0,s_1,\ldots,s_k=G_{\mathrm{goal}}$ 反复应用该条件，
有 $h(s_0)\leq \sum_{i=0}^{k-1}c(s_i,s_{i+1})+h(s_k)$。
再取最短的目标路径，就得到 $h(s)\leq h^*(s)$。
因此，对当前不重开 Closed 的 A*，证明一致性同时提供了所需的剩余代价下界。

启发函数的组合也要保持这些条件。多个一致启发函数取最大值仍然一致，
而直接相加可能重复计算同一步移动的成本，不能仅因估计更大就判为更优。
本实验的加性模式数据库需要先建立不重复的动作成本分摊，
再论证各模式距离之和的下界与一致性；具体建模放在后续的启发函数部分讨论。
\FloatBarrier

\OverviewTitle{\TemplateSubsection}{1.4 IDA*算法基本原理及其可采纳性}

A* 需要同时保留大量 Open 与 Closed 记录，在 15-Puzzle 上可能先受到存储限制。
IDA* 将按估值选择候选的思想与深度优先搜索结合：给定阈值 $B$，
仅在 $f(n)\leq B$ 的范围内继续递归，一旦超界便回溯。
这里限制的是 $g+h$，并不是统一的深度上限，因此启发估计仍然决定哪些分支能够进入本轮搜索。

\textbf{（1）阈值与一轮搜索。}
初始阈值取 $B_0=h(S)$。每次递归先计算 $f=g+h$，
若 $f>B$ 就返回该超界值；未超界时再判定是否为目标，
否则按合法动作递归并汇总各分支返回的最小超界值。
若本轮找到目标则结束；若没有找到目标而仍有超界分支，下一轮使用
\begin{equation}
  B_{k+1}=\min\{\,f(n)\mid n\text{ 在第 }k\text{ 轮被超界截断}\,\}.
  \label{eq:overview-bound}
\end{equation}
若已穷尽且没有任何超界候选，则不存在尚可继续搜索的路径。
阈值按实际超界值推进，而不是固定执行 $B\leftarrow B+1$，
这样可跳过没有候选结点进入的数值区间。

\OverviewFigure[0.78]{fig-1-5-ida-threshold}
  {IDA* 按最小超界值推进阈值的三轮搜索示意}
  {fig:overview-ida}

图~\ref{fig:overview-ida} 给出阈值 $2\to4\to6$ 的抽象例子。
第一轮只有 $f\leq2$ 的结点能够继续展开，所有被截断结点中的最小值为 4；
第二轮重新从根出发，允许进入 $f=4$ 的结点，而最小超界值变为 6；
第三轮沿当前深度优先顺序到达代价为 6 的目标后立即停止，
所以右侧某些满足阈值的分支仍呈灰色，表示尚未访问。
结点内的数字是 $f$，标注“再 3 步”的边合并了三个动作，
被省略中间结点的 $f$ 均为 6。这些值只用于解释更新规则，不是实测阈值序列。

\textbf{（2）BASIC 模式的最优性。}
以下论证针对可采纳启发函数、正代价、每轮完整搜索相应阈值范围且只做路径环剪枝的基本 IDA*。
设最优解代价为 $C^*$。由于 $B_0=h(S)\leq C^*$，
只需说明递增阈值不会跳过 $C^*$，并且到达该阈值时能找到最优解。
在一条最优路径上，任意前缀结点都满足
\begin{equation}
  f(n)=g^*(n)+h(n)\leq g^*(n)+h^*(n)=C^*.
  \label{eq:overview-ida-optimal}
\end{equation}
当 $B<C^*$ 时，目标自身的 $f=C^*$ 超界，因此该最优路径上必有第一个被截断的结点，
其返回值不超过 $C^*$；全体截断值的最小值也就不超过 $C^*$。
在本实验有限状态、单位代价且禁止路径环的条件下，候选路径及其估值有限，
阈值严格增大后最终能够到达 $C^*$。
到这一轮时，最优路径的所有前缀都通过估值检查，完整的深度优先遍历必能到达目标。

同时，算法在目标判定之前执行超界检查，因而返回目标时必有
$g(G)\leq B\leq C^*$；结合最优代价的定义，可知 $g(G)=C^*$。
这解释了为什么最后一轮即使按深度优先顺序访问后继，也不会先返回代价更大的解。
该论证只需要可采纳性，不要求每条边的 $f$ 单调。
路径环剪枝也不会删去必要的最优解，因为正代价环可以移除，得到代价更小的同终点路径。
\FloatBarrier

\textbf{（3）重复访问与存储取舍。}
IDA* 在提高阈值后需要重新展开先前访问过的路径，这是节省 Open 存储所付出的时间代价。
若把上一轮的 Closed 直接带入下一轮，某个状态虽然以前访问过，
其后继却可能因旧阈值被截断，永久关闭它就会遗漏新阈值下应当继续搜索的分支。
所以本项目不保留跨轮 Closed，当前路径上的状态在回溯时移出路径集合。

基本深度优先实现同时保存递归路径、路径状态集合及各层尚未处理的后继，
搜索存储可记为 $O(bd)$，其中 $b$ 为最大分支数、$d$ 为搜索深度；
N-Puzzle 的 $b\leq4$，因此在固定棋盘规模下随深度线性增长。
这一结论仅指 BASIC 的路径相关搜索存储，不包含模式数据库、额外置换表及其他输入存储。
此外，当前实现还记录各轮阈值，设迭代次数为 $I$，这份诊断序列额外占用 $O(I)$ 空间。
A* 则要保留大量已发现状态，存储规模随 Open 与 Closed 的总记录数增长。

本项目另有 TT 和 ZOBRIST 模式，用本轮记录的较小 $g$ 减少重复访问，
每次增加阈值前都会清空相应表。Zobrist 值用于选择桶，
只有完整状态相等时才根据 $g$ 作剪枝决定，避免把哈希碰撞直接当成同一棋盘。
这些模式需要额外表空间，其剪枝正确性还受状态记录及路径上下文处理方式约束，
不能将 BASIC 的上述证明原样当作任意置换表策略的证明。
小规模最短路径与碰撞回归可核验实现，但也不能替代对任意状态图的普遍论证。

A* 与 IDA* 的取舍因此同时涉及时间和空间：A* 保留搜索前沿，减少跨轮重新搜索，
却可能消耗较多内存；IDA* BASIC 的路径存储较小，但可能反复展开同一批结点。
启发值更接近真实代价时，两者通常都有机会缩小搜索范围，
最终效果仍受估值成本、重复状态比例、阈值轮数和可用内存影响。
后续实验应在同题、同环境和一致计时口径下比较，不能预先认定 IDA* 在所有实例上都更快。

本文第四章的主要对照实验及困难实例求解采用 IDA* 的 BASIC 模式，
开启快速实现（\texttt{wattai.fast=true}），不启用置换表。
TT 与 ZOBRIST 属于可选优化模式，其实验结果应结合具体配置分别讨论，
不能直接归入 BASIC 模式的结果。

\endgroup
\clearpage
